实迹具有循环性，且 $\operatorname{Re}\operatorname{tr}(A^2)$ 是各矩阵元模长的平方和，故此配对在 Hermitian 矩阵上是实的、对称的且非退化。特别地，$A:(vv^*)=v^*Av$。四元数 Hermitian 矩阵有标准正交特征基：在单位球面上取其二次型的最大值，并用一阶变分条件得到一个特征向量；随后在其正交补上重复，Hermitian 对称性保证该正交补不变。Newton 恒等式将行列式表示成幂和 $\operatorname{Re}\operatorname{tr}(A^j)$（$1\le j\le m$）的多项式，由此得到行列式与余子式矩阵的次数。这是 \cite{r-2}, Section~1.1 所用的 Moore 行列式约定。

在 $\mathbb R^{4m}$ 上的实表示中，每个特征值都出现四次。当 $A>0$ 时，$\log\det A$ 等于实对数行列式的四分之一。它沿方向 $B$ 的一阶导数为 $\operatorname{Re}\operatorname{tr}(A^{-1}B)$，沿直线的二阶导数为 $-\operatorname{Re}\operatorname{tr}((A^{-1/2}BA^{-1/2})^2)\le0$。对于 $B\ge0$，一阶导数非负。这些恒等式给出正定情形的余子式公式、单调性和凹性。在对角基下，余子式矩阵的特征值是其余 $m-1$ 个特征值的乘积：先对正定 $A$ 得到结论，再由多项式性推广到所有对角 $A$。由酉不变性得到关于 $A\ge0$ 及其在奇异情形下的像的断言。

在向量 $(b,w)$ 处，$M-\sigma v_0v_0^*$ 的二次型通过配平方可写为 $(w+D^{-1}cb)^*D(w+D^{-1}cb)$。若 $\sigma>0$，则在实表示中使用实 Schur 行列式公式，得到 $(\det M)^4=(\det D)^4\sigma^4$，再由正定性取正的四次方根。$\det M=(\det D)\sigma$ 的两边都是 $\ell$ 的多项式，故该结论推广到所有 $\sigma$。为证秩一公式，将单位向量 $v$ 扩充为四元数标准正交基，并用其酉矩阵 $U$，于是 $U^*v=e_1$。将命题~\ref{ext:alesker-determinant} 应用于 $U^*AU$，取 $T=\operatorname{diag}(t,0,\ldots,0)$。对角展开中只有原行列式项，以及 $t$ 乘以 $\{2,\ldots,m\}$ 上的主子式，后者为 $\det(A|_{v^\perp})$。于是任意 Hermitian 矩阵 $A$ 均满足摄动恒等式。在 $t=0$ 处求导，并采用实迹配对，得 $v^* C(A)v=\det(A|_{v^\perp})$。
